\documentclass{article}
\usepackage[T1]{fontenc}
\usepackage{lmodern}
\usepackage{graphicx}
\usepackage{amsmath,amssymb}
\usepackage[a4paper,margin=2cm]{geometry}
\usepackage[absolute,overlay]{textpos}
\setlength{\TPHorizModule}{1mm}
\setlength{\TPVertModule}{1mm}
\usepackage[indonesian]{babel}
\usepackage{hyperref}
\setlength{\emergencystretch}{2em}
% unit: Halaman 10

\begin{document}

% === Halaman 10 (python/native) ===

10

Dekripsi

1.  Hitung kembali blok plainteks m dari cipherteks c menggunakan kunci 

privat d dengan persamaan 

                   m = cd mod n,

       Jika cipherteks adalah blok c1, c2,  c3, …, dekripsi setiap blok menjadi

                   mi = ci

d mod n

2. Postprocessing: Kodekan kembali m atau mi menjadi bentuk pesan 

semula

\end{document}
